% figs.tex —— 由 ch5/make_figs.py 產生，請勿手動修改。
% 重新產生：python ch5/make_figs.py（第 5 章，動畫 W501–W507）

\newcommand{\FiveTabBiased}{%
\begin{tabular}[t]{@{}lccc@{}}
\toprule
$N=10$ & $x=0$ & $x=2$ & $x=4$\\
\midrule
精確 (5.1) & 0.2007 & 0.2508 & 0.2150\\
$2\times$(5.7) & 0.2091 & 0.2575 & 0.2091\\
\bottomrule
\end{tabular}
\qquad
\begin{tabular}[t]{@{}lccc@{}}
\toprule
Fig.~5.1 & $N=10$ & $N=20$ & $N=40$\\
\midrule
平均 $(2p-1)N$ & 2 & 4 & 8\\
變異數 $4Np(1-p)$ & 9.6 & 19.2 & 38.4\\
峰值（精確） & 0.2508 & 0.1797 & 0.1279\\
峰值（$2\times$高斯） & 0.2575 & 0.1821 & 0.1288\\
\bottomrule
\end{tabular}}

\newcommand{\FiveTabLR}{%
{\small\setlength{\tabcolsep}{3pt}\begin{tabular}{@{}cccccccccc@{}}
\toprule
$N\backslash x$ & $-4$ & $-3$ & $-2$ & $-1$ & $0$ & $1$ & $2$ & $3$ & $4$\\
\midrule
0 &  &  &  &  & $(\frac{1}{2},\frac{1}{2})$ &  &  &  & \\
1 &  &  &  & $(\frac{1}{2},0)$ &  & $(0,\frac{1}{2})$ &  &  & \\
2 &  &  & $(\frac{3}{8},0)$ &  & $(\frac{1}{8},\frac{1}{8})$ &  & $(0,\frac{3}{8})$ &  & \\
3 &  & $(\frac{9}{32},0)$ &  & $(\frac{1}{8},\frac{3}{32})$ &  & $(\frac{3}{32},\frac{1}{8})$ &  & $(0,\frac{9}{32})$ & \\
4 & $(\frac{27}{128},0)$ &  & $(\frac{15}{128},\frac{9}{128})$ &  & $(\frac{13}{128},\frac{13}{128})$ &  & $(\frac{9}{128},\frac{15}{128})$ &  & $(0,\frac{27}{128})$\\
\bottomrule
\end{tabular}}}

\newcommand{\FiveTabPVar}{%
\begin{tabular}[t]{@{}cccccc@{}}
\toprule
$Q$ & $N=1$ & $2$ & $3$ & $4$ & $n=\frac{Q}{1-Q}$\\
\midrule
$\frac{3}{4}$ & $1$ & $3$ & $\frac{11}{2}$ & $\frac{33}{4}$ & $3$\\
$\frac{1}{2}$ & $1$ & $2$ & $3$ & $4$ & $1$\\
$\frac{1}{4}$ & $1$ & $1$ & $\frac{3}{2}$ & $\frac{7}{4}$ & $\frac{1}{3}$\\
\bottomrule
\end{tabular}
\qquad
\begin{tabular}[t]{@{}lccc@{}}
\toprule
變異數 & $N=10$ & $N=100$ & $N=1000$\\
\midrule
$Q=\tfrac34$ 精確 & 26.00 & 296.00 & 2996.00\\
\quad $NQ/(1-Q)$ & 30 & 300 & 3000\\
$Q=\tfrac14$ 精確 & 3.777 & 33.778 & 333.778\\
\quad $NQ/(1-Q)$ & 3.333 & 33.333 & 333.333\\
\bottomrule
\end{tabular}}

\newcommand{\FiveFigTele}{%
\begin{tikzpicture}[x=1.3cm,y=7.5cm,>=Stealth]
\draw[->] (0,0) -- (9.5,0) node[right,font=\scriptsize] {$x$};
\draw[->] (0,0) -- (0,0.46) node[above,font=\scriptsize] {$\Sigma(x,t)$};
\foreach \x in {0,...,9} {\draw (\x,0) -- (\x,-0.008); \node[below,font=\scriptsize] at (\x,-0.006) {$\x$};}
\foreach \y in {0.1,0.2,0.3,0.4} {\draw (0,\y) -- (-0.08,\y); \node[left,font=\scriptsize] at (-0.08,\y) {$\y$};}
\draw[gray,thick,dashed] (0.000,0.14105) -- (0.150,0.14085) -- (0.300,0.14026) -- (0.450,0.13927) -- (0.600,0.13791) -- (0.750,0.13617) -- (0.900,0.13408) -- (1.050,0.13166) -- (1.200,0.12891) -- (1.350,0.12586) -- (1.500,0.12254) -- (1.650,0.11898) -- (1.800,0.11519) -- (1.950,0.11121) -- (2.100,0.10707) -- (2.250,0.10279) -- (2.400,0.09841) -- (2.550,0.09394) -- (2.700,0.08943) -- (2.850,0.08490) -- (3.000,0.08037) -- (3.150,0.07586) -- (3.300,0.07141) -- (3.450,0.06703) -- (3.600,0.06275) -- (3.750,0.05857) -- (3.900,0.05451) -- (4.050,0.05060) -- (4.200,0.04683) -- (4.350,0.04323) -- (4.500,0.03978) -- (4.650,0.03651) -- (4.800,0.03342) -- (4.950,0.03050) -- (5.100,0.02776) -- (5.250,0.02519) -- (5.400,0.02280) -- (5.550,0.02057) -- (5.700,0.01851) -- (5.850,0.01661) -- (6.000,0.01487) -- (6.150,0.01327) -- (6.300,0.01180) -- (6.450,0.01047) -- (6.600,0.00927) -- (6.750,0.00818) -- (6.900,0.00720) -- (7.050,0.00631) -- (7.200,0.00552) -- (7.350,0.00482) -- (7.500,0.00419) -- (7.650,0.00364) -- (7.800,0.00315) -- (7.950,0.00272) -- (8.100,0.00234) -- (8.250,0.00200) -- (8.400,0.00171) -- (8.550,0.00146) -- (8.700,0.00124) -- (8.850,0.00106) -- (9.000,0.00089);
\draw[cBar,thick] (0.0000,0.33684) -- (0.0145,0.33682) -- (0.0290,0.33678) -- (0.0435,0.33671) -- (0.0580,0.33662) -- (0.0725,0.33650) -- (0.0870,0.33635) -- (0.1014,0.33617) -- (0.1159,0.33597) -- (0.1304,0.33574) -- (0.1449,0.33548) -- (0.1594,0.33520) -- (0.1739,0.33489) -- (0.1884,0.33455) -- (0.2029,0.33419) -- (0.2174,0.33380) -- (0.2319,0.33338) -- (0.2464,0.33294) -- (0.2609,0.33246) -- (0.2754,0.33197) -- (0.2899,0.33144) -- (0.3043,0.33090) -- (0.3188,0.33032) -- (0.3333,0.32972) -- (0.3478,0.32909) -- (0.3623,0.32844) -- (0.3768,0.32776) -- (0.3913,0.32705) -- (0.4058,0.32632) -- (0.4203,0.32556) -- (0.4348,0.32478) -- (0.4493,0.32397) -- (0.4638,0.32314) -- (0.4783,0.32228) -- (0.4928,0.32140) -- (0.5072,0.32049) -- (0.5217,0.31955) -- (0.5362,0.31860) -- (0.5507,0.31762) -- (0.5652,0.31661) -- (0.5797,0.31558) -- (0.5942,0.31452) -- (0.6087,0.31344) -- (0.6232,0.31234) -- (0.6377,0.31122) -- (0.6522,0.31007) -- (0.6667,0.30889) -- (0.6812,0.30770) -- (0.6957,0.30648) -- (0.7101,0.30524) -- (0.7246,0.30397) -- (0.7391,0.30269) -- (0.7536,0.30138) -- (0.7681,0.30005) -- (0.7826,0.29869) -- (0.7971,0.29732) -- (0.8116,0.29592) -- (0.8261,0.29451) -- (0.8406,0.29307) -- (0.8551,0.29161) -- (0.8696,0.29013) -- (0.8841,0.28863) -- (0.8986,0.28711) -- (0.9130,0.28556) -- (0.9275,0.28400) -- (0.9420,0.28242) -- (0.9565,0.28082) -- (0.9710,0.27920) -- (0.9855,0.27757) -- (1.0000,0.27591) -- (1,0);
\draw[->,cBar,very thick] (1,0) -- (1,0.355);
\node[above,font=\scriptsize,cBar,align=center,inner sep=1pt] at (1,0.355) {$t=1$\\$0.1839$};
\draw[cWalk,thick] (0.0000,0.21991) -- (0.0435,0.21986) -- (0.0870,0.21972) -- (0.1304,0.21948) -- (0.1739,0.21914) -- (0.2174,0.21870) -- (0.2609,0.21817) -- (0.3043,0.21754) -- (0.3478,0.21682) -- (0.3913,0.21600) -- (0.4348,0.21509) -- (0.4783,0.21409) -- (0.5217,0.21299) -- (0.5652,0.21181) -- (0.6087,0.21054) -- (0.6522,0.20917) -- (0.6957,0.20772) -- (0.7391,0.20619) -- (0.7826,0.20457) -- (0.8261,0.20287) -- (0.8696,0.20109) -- (0.9130,0.19923) -- (0.9565,0.19729) -- (1.0000,0.19528) -- (1.0435,0.19319) -- (1.0870,0.19104) -- (1.1304,0.18881) -- (1.1739,0.18651) -- (1.2174,0.18415) -- (1.2609,0.18173) -- (1.3043,0.17925) -- (1.3478,0.17671) -- (1.3913,0.17411) -- (1.4348,0.17146) -- (1.4783,0.16876) -- (1.5217,0.16600) -- (1.5652,0.16321) -- (1.6087,0.16037) -- (1.6522,0.15749) -- (1.6957,0.15457) -- (1.7391,0.15161) -- (1.7826,0.14862) -- (1.8261,0.14560) -- (1.8696,0.14256) -- (1.9130,0.13948) -- (1.9565,0.13639) -- (2.0000,0.13328) -- (2.0435,0.13014) -- (2.0870,0.12700) -- (2.1304,0.12384) -- (2.1739,0.12068) -- (2.2174,0.11751) -- (2.2609,0.11433) -- (2.3043,0.11115) -- (2.3478,0.10798) -- (2.3913,0.10481) -- (2.4348,0.10165) -- (2.4783,0.09849) -- (2.5217,0.09535) -- (2.5652,0.09222) -- (2.6087,0.08911) -- (2.6522,0.08602) -- (2.6957,0.08295) -- (2.7391,0.07990) -- (2.7826,0.07688) -- (2.8261,0.07388) -- (2.8696,0.07092) -- (2.9130,0.06799) -- (2.9565,0.06509) -- (3.0000,0.06223) -- (3,0);
\draw[->,cWalk,very thick] (3,0) -- (3,0.355);
\node[above,font=\scriptsize,cWalk,align=center,inner sep=1pt] at (3,0.355) {$t=3$\\$0.02489$};
\draw[cExact,thick] (0.0000,0.17376) -- (0.0725,0.17368) -- (0.1449,0.17346) -- (0.2174,0.17309) -- (0.2899,0.17258) -- (0.3623,0.17191) -- (0.4348,0.17111) -- (0.5072,0.17016) -- (0.5797,0.16907) -- (0.6522,0.16785) -- (0.7246,0.16648) -- (0.7971,0.16499) -- (0.8696,0.16336) -- (0.9420,0.16161) -- (1.0145,0.15974) -- (1.0870,0.15774) -- (1.1594,0.15563) -- (1.2319,0.15342) -- (1.3043,0.15109) -- (1.3768,0.14866) -- (1.4493,0.14614) -- (1.5217,0.14352) -- (1.5942,0.14082) -- (1.6667,0.13804) -- (1.7391,0.13518) -- (1.8116,0.13225) -- (1.8841,0.12926) -- (1.9565,0.12620) -- (2.0290,0.12310) -- (2.1014,0.11994) -- (2.1739,0.11675) -- (2.2464,0.11351) -- (2.3188,0.11025) -- (2.3913,0.10697) -- (2.4638,0.10366) -- (2.5362,0.10034) -- (2.6087,0.09702) -- (2.6812,0.09369) -- (2.7536,0.09037) -- (2.8261,0.08706) -- (2.8986,0.08376) -- (2.9710,0.08048) -- (3.0435,0.07723) -- (3.1159,0.07400) -- (3.1884,0.07081) -- (3.2609,0.06766) -- (3.3333,0.06455) -- (3.4058,0.06149) -- (3.4783,0.05848) -- (3.5507,0.05553) -- (3.6232,0.05263) -- (3.6957,0.04980) -- (3.7681,0.04703) -- (3.8406,0.04433) -- (3.9130,0.04170) -- (3.9855,0.03914) -- (4.0580,0.03666) -- (4.1304,0.03426) -- (4.2029,0.03193) -- (4.2754,0.02968) -- (4.3478,0.02752) -- (4.4203,0.02544) -- (4.4928,0.02344) -- (4.5652,0.02152) -- (4.6377,0.01969) -- (4.7101,0.01794) -- (4.7826,0.01628) -- (4.8551,0.01470) -- (4.9275,0.01321) -- (5.0000,0.01179) -- (5,0);
\draw[->,cExact,very thick] (5,0) -- (5,0.355);
\node[above,font=\scriptsize,cExact,align=center,inner sep=1pt] at (5,0.355) {$t=5$\\$0.003369$};
\draw[violet!75!black,thick] (0.0000,0.13879) -- (0.1159,0.13868) -- (0.2319,0.13838) -- (0.3478,0.13787) -- (0.4638,0.13716) -- (0.5797,0.13625) -- (0.6957,0.13514) -- (0.8116,0.13384) -- (0.9275,0.13236) -- (1.0435,0.13070) -- (1.1594,0.12886) -- (1.2754,0.12686) -- (1.3913,0.12470) -- (1.5072,0.12238) -- (1.6232,0.11993) -- (1.7391,0.11734) -- (1.8551,0.11463) -- (1.9710,0.11180) -- (2.0870,0.10887) -- (2.2029,0.10584) -- (2.3188,0.10273) -- (2.4348,0.09955) -- (2.5507,0.09630) -- (2.6667,0.09301) -- (2.7826,0.08967) -- (2.8986,0.08630) -- (3.0145,0.08292) -- (3.1304,0.07952) -- (3.2464,0.07613) -- (3.3623,0.07274) -- (3.4783,0.06938) -- (3.5942,0.06604) -- (3.7101,0.06274) -- (3.8261,0.05949) -- (3.9420,0.05629) -- (4.0580,0.05316) -- (4.1739,0.05009) -- (4.2899,0.04709) -- (4.4058,0.04418) -- (4.5217,0.04135) -- (4.6377,0.03861) -- (4.7536,0.03596) -- (4.8696,0.03341) -- (4.9855,0.03096) -- (5.1014,0.02862) -- (5.2174,0.02637) -- (5.3333,0.02424) -- (5.4493,0.02221) -- (5.5652,0.02028) -- (5.6812,0.01846) -- (5.7971,0.01675) -- (5.9130,0.01514) -- (6.0290,0.01364) -- (6.1449,0.01223) -- (6.2609,0.01093) -- (6.3768,0.00972) -- (6.4928,0.00860) -- (6.6087,0.00757) -- (6.7246,0.00663) -- (6.8406,0.00577) -- (6.9565,0.00499) -- (7.0725,0.00428) -- (7.1884,0.00364) -- (7.3043,0.00308) -- (7.4203,0.00257) -- (7.5362,0.00212) -- (7.6522,0.00173) -- (7.7681,0.00139) -- (7.8841,0.00109) -- (8.0000,0.00084) -- (8,0);
\draw[->,violet!75!black,very thick] (8,0) -- (8,0.355);
\node[above,font=\scriptsize,violet!75!black,align=center,inner sep=1pt] at (8,0.355) {$t=8$\\$0.0001677$};
\end{tikzpicture}}

\newcommand{\FiveTabTele}{%
\begin{tabular}{@{}lcccc@{}}
\toprule
$v=\gamma=1$ & $t=1$ & $t=3$ & $t=5$ & $t=8$\\
\midrule
$\Sigma(0,t)$（(5.61)） & 0.3368 & 0.2199 & 0.1738 & 0.1388\\
高斯 $1/\sqrt{4\pi Dt}$（$D=\frac12$） & 0.3989 & 0.2303 & 0.1784 & 0.1410\\
前緣後方 $\Sigma(vt^-,t)=\frac{e^{-t}}2(1+\frac t2)$ & 0.2759 & 0.06223 & 0.01179 & 0.0008387\\
每個尖峰 $\frac12e^{-t}$ & 0.1839 & 0.02489 & 0.003369 & 0.0001677\\
連續部分總量 $1-e^{-t}$ & 0.6321 & 0.9502 & 0.9933 & 0.9997\\
變異數 $t-\frac{1-e^{-2t}}2$ & 0.5677 & 2.5012 & 4.5000 & 7.5000\\
\bottomrule
\end{tabular}}

\newcommand{\FiveFigRT}{%
\begin{tikzpicture}[x=1.3cm,y=1.9cm,>=Stealth]
\draw[->] (0,0) -- (8.5,0) node[right,font=\scriptsize] {$r$};
\draw[->] (0,0) -- (0,1.75) node[above,font=\scriptsize] {$P_3(r,t)=4\pi r^2\rho$};
\foreach \x in {0,...,8} {\draw (\x,0) -- (\x,-0.03); \node[below,font=\scriptsize] at (\x,-0.02) {$\x$};}
\foreach \y in {0.5,1,1.5} {\draw (0,\y) -- (-0.08,\y); \node[left,font=\scriptsize] at (-0.08,\y) {$\y$};}
\draw[gray!30] (0,0.5) -- (8.2,0.5) (0,1) -- (8.2,1) (0,1.5) -- (8.2,1.5);
\begin{scope}
\clip (0,0) rectangle (8.3,1.5);
\draw[cBar,thick] (0.0020,0.0000) -- (0.0151,0.0006) -- (0.0282,0.0019) -- (0.0413,0.0041) -- (0.0543,0.0072) -- (0.0674,0.0110) -- (0.0805,0.0157) -- (0.0935,0.0212) -- (0.1066,0.0275) -- (0.1196,0.0346) -- (0.1325,0.0424) -- (0.1455,0.0510) -- (0.1585,0.0603) -- (0.1714,0.0704) -- (0.1843,0.0812) -- (0.1971,0.0927) -- (0.2099,0.1049) -- (0.2227,0.1177) -- (0.2355,0.1311) -- (0.2482,0.1452) -- (0.2608,0.1598) -- (0.2734,0.1750) -- (0.2860,0.1908) -- (0.2985,0.2070) -- (0.3110,0.2238) -- (0.3234,0.2410) -- (0.3358,0.2587) -- (0.3481,0.2767) -- (0.3603,0.2952) -- (0.3725,0.3140) -- (0.3846,0.3331) -- (0.3967,0.3526) -- (0.4087,0.3723) -- (0.4206,0.3922) -- (0.4324,0.4124) -- (0.4442,0.4328) -- (0.4559,0.4534) -- (0.4675,0.4741) -- (0.4791,0.4949) -- (0.4905,0.5159) -- (0.5019,0.5369) -- (0.5132,0.5579) -- (0.5244,0.5790) -- (0.5355,0.6001) -- (0.5465,0.6212) -- (0.5574,0.6422) -- (0.5682,0.6632) -- (0.5789,0.6842) -- (0.5896,0.7050) -- (0.6001,0.7257) -- (0.6105,0.7463) -- (0.6208,0.7668) -- (0.6310,0.7871) -- (0.6411,0.8073) -- (0.6511,0.8273) -- (0.6610,0.8471) -- (0.6708,0.8668) -- (0.6804,0.8862) -- (0.6900,0.9055) -- (0.6994,0.9245) -- (0.7087,0.9433) -- (0.7179,0.9620) -- (0.7269,0.9804) -- (0.7359,0.9986) -- (0.7447,1.0166) -- (0.7533,1.0343) -- (0.7619,1.0519) -- (0.7703,1.0693) -- (0.7786,1.0865) -- (0.7867,1.1035) -- (0.7947,1.1204) -- (0.8026,1.1370) -- (0.8104,1.1536) -- (0.8180,1.1700) -- (0.8254,1.1863) -- (0.8328,1.2024) -- (0.8399,1.2185) -- (0.8470,1.2346) -- (0.8539,1.2506) -- (0.8606,1.2666) -- (0.8672,1.2826) -- (0.8736,1.2986) -- (0.8799,1.3147) -- (0.8861,1.3309) -- (0.8921,1.3473) -- (0.8979,1.3638) -- (0.9036,1.3806) -- (0.9091,1.3976) -- (0.9145,1.4149) -- (0.9197,1.4326) -- (0.9248,1.4508) -- (0.9297,1.4694) -- (0.9344,1.4886) -- (0.9390,1.5084) -- (0.9434,1.5289) -- (0.9477,1.5503) -- (0.9518,1.5726) -- (0.9557,1.5959) -- (0.9595,1.6204) -- (0.9631,1.6462) -- (0.9665,1.6735) -- (0.9698,1.7024) -- (0.9729,1.7333) -- (0.9759,1.7663) -- (0.9786,1.8018) -- (0.9813,1.8402) -- (0.9837,1.8818) -- (0.9860,1.9273) -- (0.9881,1.9772) -- (0.9900,2.0324);
\draw[cWalk,thick] (0.0060,0.0000) -- (0.0453,0.0007) -- (0.0846,0.0025) -- (0.1238,0.0054) -- (0.1630,0.0093) -- (0.2023,0.0142) -- (0.2414,0.0202) -- (0.2806,0.0272) -- (0.3197,0.0351) -- (0.3587,0.0439) -- (0.3976,0.0535) -- (0.4365,0.0640) -- (0.4754,0.0753) -- (0.5141,0.0873) -- (0.5528,0.0999) -- (0.5913,0.1131) -- (0.6298,0.1269) -- (0.6681,0.1411) -- (0.7064,0.1557) -- (0.7445,0.1707) -- (0.7825,0.1860) -- (0.8203,0.2014) -- (0.8581,0.2170) -- (0.8956,0.2327) -- (0.9330,0.2484) -- (0.9703,0.2640) -- (1.0074,0.2796) -- (1.0443,0.2949) -- (1.0810,0.3101) -- (1.1176,0.3249) -- (1.1539,0.3394) -- (1.1901,0.3536) -- (1.2261,0.3673) -- (1.2618,0.3805) -- (1.2973,0.3932) -- (1.3326,0.4054) -- (1.3677,0.4170) -- (1.4026,0.4281) -- (1.4372,0.4384) -- (1.4715,0.4482) -- (1.5056,0.4573) -- (1.5395,0.4657) -- (1.5731,0.4734) -- (1.6064,0.4804) -- (1.6394,0.4868) -- (1.6722,0.4924) -- (1.7046,0.4974) -- (1.7368,0.5016) -- (1.7687,0.5052) -- (1.8003,0.5082) -- (1.8315,0.5104) -- (1.8625,0.5121) -- (1.8931,0.5131) -- (1.9234,0.5135) -- (1.9534,0.5134) -- (1.9831,0.5127) -- (2.0124,0.5115) -- (2.0413,0.5097) -- (2.0699,0.5075) -- (2.0982,0.5049) -- (2.1261,0.5018) -- (2.1536,0.4983) -- (2.1808,0.4944) -- (2.2076,0.4902) -- (2.2340,0.4857) -- (2.2600,0.4809) -- (2.2856,0.4758) -- (2.3109,0.4705) -- (2.3357,0.4650) -- (2.3602,0.4593) -- (2.3842,0.4534) -- (2.4079,0.4474) -- (2.4311,0.4413) -- (2.4539,0.4351) -- (2.4763,0.4289) -- (2.4983,0.4226) -- (2.5198,0.4162) -- (2.5409,0.4099) -- (2.5616,0.4036) -- (2.5818,0.3974) -- (2.6016,0.3911) -- (2.6209,0.3850) -- (2.6398,0.3789) -- (2.6582,0.3730) -- (2.6762,0.3672) -- (2.6937,0.3615) -- (2.7108,0.3559) -- (2.7274,0.3505) -- (2.7435,0.3453) -- (2.7592,0.3402) -- (2.7744,0.3354) -- (2.7891,0.3308) -- (2.8033,0.3264) -- (2.8170,0.3222) -- (2.8303,0.3183) -- (2.8431,0.3146) -- (2.8554,0.3113) -- (2.8672,0.3082) -- (2.8785,0.3055) -- (2.8893,0.3031) -- (2.8996,0.3011) -- (2.9095,0.2994) -- (2.9188,0.2982) -- (2.9276,0.2974) -- (2.9359,0.2972) -- (2.9438,0.2975) -- (2.9511,0.2983) -- (2.9579,0.2999) -- (2.9642,0.3023) -- (2.9700,0.3055);
\draw[cExact,thick] (0.0100,0.0000) -- (0.0755,0.0009) -- (0.1409,0.0030) -- (0.2064,0.0064) -- (0.2717,0.0110) -- (0.3371,0.0168) -- (0.4024,0.0238) -- (0.4676,0.0319) -- (0.5328,0.0410) -- (0.5978,0.0511) -- (0.6627,0.0621) -- (0.7276,0.0738) -- (0.7923,0.0863) -- (0.8569,0.0993) -- (0.9213,0.1129) -- (0.9856,0.1269) -- (1.0497,0.1411) -- (1.1136,0.1556) -- (1.1773,0.1702) -- (1.2408,0.1848) -- (1.3042,0.1993) -- (1.3672,0.2135) -- (1.4301,0.2275) -- (1.4927,0.2412) -- (1.5551,0.2543) -- (1.6172,0.2670) -- (1.6790,0.2790) -- (1.7405,0.2904) -- (1.8017,0.3010) -- (1.8626,0.3110) -- (1.9232,0.3201) -- (1.9835,0.3283) -- (2.0434,0.3357) -- (2.1030,0.3422) -- (2.1622,0.3479) -- (2.2211,0.3526) -- (2.2795,0.3564) -- (2.3376,0.3594) -- (2.3953,0.3614) -- (2.4525,0.3626) -- (2.5094,0.3630) -- (2.5658,0.3625) -- (2.6218,0.3613) -- (2.6773,0.3594) -- (2.7323,0.3567) -- (2.7869,0.3534) -- (2.8410,0.3494) -- (2.8947,0.3449) -- (2.9478,0.3399) -- (3.0004,0.3343) -- (3.0526,0.3283) -- (3.1041,0.3220) -- (3.1552,0.3152) -- (3.2057,0.3082) -- (3.2557,0.3009) -- (3.3051,0.2934) -- (3.3539,0.2857) -- (3.4022,0.2778) -- (3.4499,0.2699) -- (3.4970,0.2618) -- (3.5435,0.2538) -- (3.5894,0.2457) -- (3.6346,0.2377) -- (3.6793,0.2296) -- (3.7233,0.2217) -- (3.7667,0.2139) -- (3.8094,0.2061) -- (3.8515,0.1985) -- (3.8929,0.1911) -- (3.9336,0.1838) -- (3.9737,0.1766) -- (4.0131,0.1697) -- (4.0518,0.1629) -- (4.0898,0.1564) -- (4.1271,0.1500) -- (4.1638,0.1439) -- (4.1996,0.1380) -- (4.2348,0.1322) -- (4.2693,0.1267) -- (4.3030,0.1214) -- (4.3359,0.1163) -- (4.3682,0.1115) -- (4.3997,0.1068) -- (4.4304,0.1023) -- (4.4603,0.0981) -- (4.4896,0.0940) -- (4.5180,0.0901) -- (4.5456,0.0865) -- (4.5725,0.0830) -- (4.5986,0.0796) -- (4.6239,0.0765) -- (4.6484,0.0735) -- (4.6722,0.0708) -- (4.6951,0.0681) -- (4.7172,0.0656) -- (4.7385,0.0633) -- (4.7590,0.0612) -- (4.7786,0.0591) -- (4.7975,0.0573) -- (4.8155,0.0555) -- (4.8327,0.0539) -- (4.8491,0.0525) -- (4.8647,0.0511) -- (4.8794,0.0499) -- (4.8932,0.0489) -- (4.9063,0.0480) -- (4.9185,0.0472) -- (4.9298,0.0466) -- (4.9403,0.0462) -- (4.9500,0.0459);
\draw[violet!75!black,thick] (0.0160,0.0000) -- (0.1207,0.0010) -- (0.2255,0.0036) -- (0.3302,0.0076) -- (0.4348,0.0132) -- (0.5394,0.0201) -- (0.6438,0.0283) -- (0.7482,0.0377) -- (0.8524,0.0481) -- (0.9565,0.0595) -- (1.0604,0.0718) -- (1.1641,0.0846) -- (1.2677,0.0980) -- (1.3710,0.1118) -- (1.4741,0.1257) -- (1.5769,0.1397) -- (1.6795,0.1535) -- (1.7817,0.1671) -- (1.8837,0.1803) -- (1.9853,0.1930) -- (2.0866,0.2050) -- (2.1876,0.2163) -- (2.2882,0.2267) -- (2.3883,0.2362) -- (2.4881,0.2447) -- (2.5874,0.2521) -- (2.6863,0.2585) -- (2.7848,0.2637) -- (2.8827,0.2678) -- (2.9802,0.2708) -- (3.0772,0.2727) -- (3.1736,0.2735) -- (3.2695,0.2732) -- (3.3648,0.2720) -- (3.4595,0.2697) -- (3.5537,0.2666) -- (3.6472,0.2627) -- (3.7401,0.2580) -- (3.8324,0.2525) -- (3.9240,0.2465) -- (4.0150,0.2399) -- (4.1052,0.2329) -- (4.1948,0.2254) -- (4.2836,0.2176) -- (4.3717,0.2096) -- (4.4591,0.2014) -- (4.5457,0.1930) -- (4.6315,0.1845) -- (4.7165,0.1761) -- (4.8007,0.1676) -- (4.8841,0.1593) -- (4.9666,0.1510) -- (5.0483,0.1429) -- (5.1291,0.1350) -- (5.2091,0.1273) -- (5.2881,0.1199) -- (5.3663,0.1126) -- (5.4435,0.1057) -- (5.5198,0.0990) -- (5.5952,0.0926) -- (5.6696,0.0865) -- (5.7430,0.0807) -- (5.8154,0.0752) -- (5.8868,0.0700) -- (5.9573,0.0651) -- (6.0267,0.0604) -- (6.0950,0.0560) -- (6.1624,0.0519) -- (6.2286,0.0481) -- (6.2938,0.0445) -- (6.3580,0.0411) -- (6.4210,0.0380) -- (6.4829,0.0351) -- (6.5437,0.0324) -- (6.6034,0.0299) -- (6.6620,0.0275) -- (6.7194,0.0254) -- (6.7757,0.0234) -- (6.8308,0.0215) -- (6.8847,0.0198) -- (6.9375,0.0183) -- (6.9891,0.0168) -- (7.0394,0.0155) -- (7.0886,0.0143) -- (7.1366,0.0132) -- (7.1833,0.0122) -- (7.2288,0.0112) -- (7.2730,0.0104) -- (7.3160,0.0096) -- (7.3578,0.0088) -- (7.3983,0.0082) -- (7.4375,0.0076) -- (7.4754,0.0070) -- (7.5121,0.0065) -- (7.5475,0.0061) -- (7.5816,0.0056) -- (7.6144,0.0053) -- (7.6458,0.0049) -- (7.6760,0.0046) -- (7.7049,0.0043) -- (7.7324,0.0041) -- (7.7586,0.0038) -- (7.7835,0.0036) -- (7.8070,0.0034) -- (7.8292,0.0032) -- (7.8501,0.0031) -- (7.8696,0.0030) -- (7.8877,0.0029) -- (7.9045,0.0028) -- (7.9200,0.0027);
\end{scope}
\draw[->,cBar,very thick] (1,0) -- (1,1.3);
\node[above,font=\scriptsize,cBar,align=center,inner sep=1pt,fill=white] at (1,1.3) {$t=1$\\$0.3679$};
\draw[->,cWalk,very thick] (3,0) -- (3,1.3);
\node[above,font=\scriptsize,cWalk,align=center,inner sep=1pt,fill=white] at (3,1.3) {$t=3$\\$0.04979$};
\draw[->,cExact,very thick] (5,0) -- (5,1.3);
\node[above,font=\scriptsize,cExact,align=center,inner sep=1pt,fill=white] at (5,1.3) {$t=5$\\$0.006738$};
\draw[->,violet!75!black,very thick] (8,0) -- (8,1.3);
\node[above,font=\scriptsize,violet!75!black,align=center,inner sep=1pt,fill=white] at (8,1.3) {$t=8$\\$0.0003355$};
\end{tikzpicture}}

\newcommand{\FiveTabRT}{%
\begin{tabular}[t]{@{}lccccc@{}}
\toprule
$P_3$ & $r=\frac t4$ & $\frac t2$ & $\frac{3t}4$ & $0.9t$ & 球殼\\
\midrule
$t=1$ & 0.1472 & 0.5334 & 1.0275 & 1.3699 & 0.3679\\
$t=3$ & 0.1729 & 0.4558 & 0.4828 & 0.3594 & 0.04979\\
$t=5$ & 0.1869 & 0.3630 & 0.2169 & 0.0926 & 0.006738\\
$t=8$ & 0.1948 & 0.2411 & 0.0622 & 0.0118 & 0.0003355\\
\bottomrule
\end{tabular}}

\newcommand{\FiveTabDims}{%
\begin{tabular}[t]{@{}lccc@{}}
\toprule
$t=3$ & 一維 & 二維 & 三維\\
\midrule
沒轉過：$e^{-3}$ & 0.0498 & 0.0498 & 0.0498\\
沒轉過：模擬 & 0.0488 & 0.0493 & 0.0497\\
$r<1.5$：公式 & 0.6051 & 0.3310 & 0.2881\\
$r<1.5$：模擬 & 0.6060 & 0.3335 & 0.2880\\
前緣 & 0.06223 & $1/\sqrt{\ }$ 發散 & $\ln$ 發散\\
\bottomrule
\end{tabular}}

\newcommand{\FiveTabCTRW}{%
{\small\setlength{\tabcolsep}{4pt}\begin{tabular}{@{}lcccccccc@{}}
\toprule
$t=3$ 走了 $m$ 步 & $0$ & $1$ & $2$ & $3$ & $4$ & $5$ & $6$ & 平均\\
\midrule
固定速率 $p_0=1$ & 0.0498 & 0.1494 & 0.2240 & 0.2240 & 0.1680 & 0.1008 & 0.0504 & 3\\
兩段等待 & 0.1991 & 0.4481 & 0.2689 & 0.0720 & 0.0108 & 0.0010 & 0.0001 & 1.2506\\
\bottomrule
\end{tabular}}
\\[6pt]
\begin{tabular}{@{}lcc@{}}
\toprule
一維格子 $C=\frac12$ & 固定速率 $B(0)=1$ & 兩段等待 $B(0)=2$\\
\midrule
$P(0,10)$：精確 & 0.1278 & 0.1805\\
$P(0,10)$：高斯 $\sqrt{B(0)/2\pi t}$ & 0.1262 & 0.1784\\
變異數（精確）$t=10\to40$ & 10.00 $\to$ 40.00 & 4.75 $\to$ 19.75\\
寬度 $t=10\to40$ & 3.162 $\to$ 6.325 & 2.179 $\to$ 4.444\\
\bottomrule
\end{tabular}}

\newcommand{\FiveTabErrata}{%
\begin{tabular}{@{}lll@{}}
\toprule
位置 & 書上 & 應為\\
\midrule
p.~96，(5.3) & {\color{cWrong}$l(2p-1)\frac{\partial P}{\partial y}$} & {\color{cExact}$l(1-2p)\frac{\partial P}{\partial y}$}\\
p.~103，(5.37) & {\color{cWrong}$\lambda^2-2Q\cos ql\,\lambda+(1-2Q)=0$} & {\color{cExact}$\lambda^2-2Q\cos ql\,\lambda-(1-2Q)=0$}\\
p.~107，(5.57) 前 & {\color{cWrong}when $x<t$} & {\color{cExact}$x<vt$}\\
p.~107，(5.60)(5.61) & {\color{cWrong}$\frac12\big(2\gamma+\frac{\partial}{\partial t}\big)$} & {\color{cExact}$\frac{1}{2v}\big(2\gamma+\frac{\partial}{\partial t}\big)$}\\
p.~112，(5.83) 後 & {\color{cWrong}$b\to i\gamma$} & {\color{cExact}$b\to i\gamma/v$}\\
p.~112，(5.84) & {\color{cWrong}$2\gamma\sqrt{\pi/2}\,e^{-\gamma t}\cosh(\gamma\sqrt{v^2t^2-r^2})$} & {\color{cExact}$2\gamma\,e^{-\gamma t}\cosh(\gamma\sqrt{v^2t^2-r^2}/v)$}\\
p.~113，(5.88) & {\color{cWrong}$\frac{2\pi}{x}\delta(x-y)$} & {\color{cExact}$\frac1x\delta(x-y)$}\\
p.~113，(5.89) & {\color{cWrong}兩項前面各有 $2\pi$} & {\color{cExact}兩個 $2\pi$ 都去掉}\\
p.~115，(5.96) 前 & {\color{cWrong}$\frac1{vq}\big(\frac\pi2-i\ln((1+x)(1-x))\big)$} & {\color{cExact}$\frac1{vq}\big(\frac\pi2-\frac i2\ln\frac{1+x}{1-x}\big)$}\\
p.~117，(5.106) & {\color{cWrong}$e^{-\gamma ut\cot u}$} & {\color{cExact}$e^{+\gamma ut\cot u}$}\\
p.~124，(5.133) 後 & {\color{cWrong}width proportional to $t$} & {\color{cExact}變異數 $\propto t$（寬度 $\propto\sqrt t$）}\\
\bottomrule
\end{tabular}}


\section{漫步的變化形（第 5 章）}

\begin{keybox}[先看這裡]
前四章的醉漢最「乾淨」：每一步左右各半、和過去無關、每單位時間走一步。這一章一次改一條規則，看結果怎麼變。
\textbf{偏向}（硬幣不公正 $\to$ 高斯往前走、寬度乘上 $\sqrt{4p(1-p)}$）；
\textbf{慣性}（記得上一步 $\to$ 兩個狀態、$2\times2$ 轉移矩陣，走久了步長放大成持續長度 $l_p=\frac{Q}{1-Q}l$；時間變連續就是電報方程，二三維是細菌的「跑一跑、轉個彎」）；
\textbf{等待}（每一步之前先等一段隨機時間 $\to$ Laplace 轉換把卷積變乘法，$\frac{B}{1-\chi A}$）。
結論都一樣：走得夠久又是高斯，只是平均、寬度與時間尺度換了樣子。總整理在 5.5 節。
\textbf{記號}：步長 $l=1$；$p$＝向右的機率；$Q$＝和上一步同方向的機率（第一步左右各半）；$L$、$R$＝最後一步向左／向右的機率；
速率 $v$、反轉（二三維：轉向）率 $\gamma$，例子都取 $v=\gamma=1$；密度標準歸一化，徑向密度 $P_2=2\pi r\rho$、$P_3=4\pi r^2\rho$。

\textbf{全章的主例}：$p=0.6$ 的硬幣走 10 步，平均 2、變異數 9.6；$Q=\frac34$ 的醉漢走 1000 步，寬度相當於普通醉漢走 3000 步；
電報走路人 $t=3$ 時還有 $e^{-3}=0.0498$ 沒轉過身；每步前先「等公車」，$t=3$ 之前平均走 3 步或 1.2506 步。
\begin{center}
\small
\begin{tabular}{@{}lll@{\qquad}lll@{}}
\toprule
節 & 內容 & 動畫 & 節 & 內容 & 動畫\\
\midrule
5.1 & 不公正的硬幣、先推再走 & W501 & 5.2.3 & 電報方程 & W504\\
5.2.1 & 兩個狀態、轉移矩陣 & W502 & 5.2.4 & 二維、三維：轉個彎 & W505\\
5.2.2 & 持續長度、二階遞迴 & W503 & 5.3.1--5.3.2 & 等待、高斯極限、會消失 & W506--W507\\
\bottomrule
\end{tabular}
\end{center}
數字一律由 \texttt{codes/utils/walk\_var.py} 重算，圖表由 \texttt{ch5/make\_figs.py} 產生並先做數字檢查；逐式校對見 \texttt{ref/md/ch05\_formulas.md}。
連續播放：\texttt{Ch5Demo}（W501--W507）。
\end{keybox}

% =====================================================================
\subsection{有偏向的漫步}
\leavevmode\bk{pp.~95--98}\ \anim{W501}\quad 硬幣不公正，或每一步都被「風」推一下：平均位置開始移動。問題是寬度會不會跟著變。

\textbf{問}：向右的機率 $p=0.6$。醉漢走 10 步，平均到哪裡？分布比公正硬幣寬還是窄？

\textbf{直觀}：一步的位移 $\pm1$，平均 $p-(1-p)=2p-1=0.2$，變異數 $1-(2p-1)^2=4p(1-p)=0.96$（公正硬幣是 1）。步與步獨立，$N$ 步的平均與變異數都乘 $N$：
平均 $(2p-1)N=2$、變異數 $4Np(1-p)=9.6$。\textbf{揭曉}：偏向讓分布\textbf{變窄}一點——向右的步多了，結果比較「確定」。

\textbf{算到底}：$n$ 步向右、$N-n$ 步向左，終點 $x=2n-N$；也可以一步一步遞迴：
\[P(N;n,N-n)=\binom Nn p^n(1-p)^{N-n},\qquad P(N+1;x,y)=p\,P(N;x,y-l)+(1-p)\,P(N;x,y+l).\tag{5.1, 5.2}\]
$x=2$ 要 6 步向右：$P(10;2)=\binom{10}{6}0.6^6\,0.4^4=210\times0.046656\times0.0256=0.2508$；同樣 $P(10;0)=\binom{10}{5}(0.24)^5=0.2007$、$P(10;4)=0.2150$。

\textbf{高斯}：照第 1 章取傅立葉，一步貢獻 $pe^{iql}+(1-p)e^{-iql}$，取對數展開到 $q^2$：
\begin{gather*}
P(N;q)=\exp\big[N\,il(2p-1)q-2Nq^2l^2p(1-p)+O(q^3)\big]P(0;q),\tag{5.5c}\\
P(N;x)\approx\frac{1}{\sqrt{8\pi Np(1-p)l^2}}\exp\Big[-\frac{(x-(2p-1)Nl)^2}{8Nl^2p(1-p)}\Big].\tag{5.7}
\end{gather*}
只有和 $N$ 同奇偶的 $x$ 有人，所以格點機率 $\approx2\times$(5.7)：$x=0,2,4$ 是 0.2091、0.2575、0.2091（精確 0.2007、0.2508、0.2150）。
高斯對平均 2 對稱，精確分布不對稱（$x=4$ 比 $x=0$ 多，差在 $q^3$ 項）。$N$ 加大，高斯往右走、變寬（Fig.~5.1，表 \ref{five-tab-biased}）；
寬度 $\propto\sqrt{p(1-p)}$，$p\to1$ 時 $\to0$：沒有隨機性了。

\begin{table}[htbp]
\centering
\small
\FiveTabBiased
\caption{不公正的硬幣 $p=0.6$。左：$N=10$ 的精確值與 $2\times$ 高斯；右：Fig.~5.1 的三組，平均與變異數都和 $N$ 成正比，峰值約除以 $\sqrt2$}\label{five-tab-biased}
\end{table}

\textbf{擴散方程的版本}：把 (5.2) 對 $y$ 泰勒展開，
\[P(N+1;x,y)=P(N;x,y)+l(1-2p)\frac{\partial P}{\partial y}+\frac{l^2}{2}\frac{\partial^2P}{\partial y^2}+\cdots\tag{5.3}\]
\begin{errbox}[p.~96，(5.3)]
書上一階項寫成 $l(2p-1)\frac{\partial P}{\partial y}$。$pP(y-l)+(1-p)P(y+l)$ 的一階項是 $-pl\,\partial_yP+(1-p)l\,\partial_yP=l(1-2p)\,\partial_yP$：往右偏、峰往右移，前面要是負號。
取傅立葉（$\int e^{iqy}\partial_yP\,dy=-iqP$）之後才得到 (5.4) 的 $+il(2p-1)q$——(5.4) 是對的。
\end{errbox}

\textbf{練習 5.1 的陷阱}：把 $P(N+1)-P(N)$ 直接換成 $\partial P/\partial N$，得擴散係數 $D=\frac{l^2}2$、變異數 $Nl^2=10$、寬度 $\sqrt{10}=3.162$——和偏向無關，錯了。
多保留一項 $\frac12\partial_N^2P$，再用一階的 $\partial_NP\approx-v\,\partial_yP$（$v=(2p-1)l$）把它換成 $\frac{v^2}2\partial_y^2P$，得
$D=\frac{l^2}{2}-\frac{v^2}{2}=2p(1-p)\,l^2$。$p=0.6$：$D=0.48$、變異數 $2DN=9.6$、寬度 3.098，和精確解一致。兩者的寬度比 $\sqrt{4p(1-p)}$ 在 $p=0.6$ 是 0.9798（差 2\%），$p=0.9$ 是 0.6（差 40\%）：
擴散方程只在 $p\approx\frac12$ 時可以直接寫。

\textbf{先走一步再亂走}（書 5.1.1 節，pp.~97--98）：另一種偏向是每一步先被固定的位移 $\vec r_B$ 推一下（像水流帶著染料），再照原來的分布亂走。一步的特徵函數只多一個相位：
\[\chi(\vec q)=\int e^{-i\vec q\cdot\vec r}P(\vec r-\vec r_B)\,d^3r=e^{-i\vec q\cdot\vec r_B}\chi_0(\vec q),\qquad
\frac{1}{(2\pi)^3}\int d^3q\,\chi(\vec q)^Ne^{i\vec q\cdot\vec r}=P(N;\vec r-N\vec r_B).\tag{5.9, 5.10}\]
整團分布平移 $N\vec r_B$，\textbf{寬度完全不變}。一維對照（$N=10$）：每步先被吹 $b=0.2$ 再擲公正硬幣，平均 2、變異數 10；不公正硬幣 $p=0.6$，平均 2、變異數 9.6。
平均一樣，但「不公正」拿掉了一點隨機性，「被吹」沒有。

% =====================================================================
\subsection{有慣性的漫步}
\leavevmode\bk{pp.~98--118}\quad 醉漢記得上一步：有 $Q$ 的機率照同一個方向走、$1-Q$ 的機率回頭（$Q=\frac12$ 就是第 1 章）。只看位置已經不是「沒有記憶」的漫步；
把「最後一步的方向」也記下來，又變回乘矩陣。

\subsubsection{兩個狀態與轉移矩陣}
\leavevmode\bk{pp.~98--100}\ \anim{W502}\par\nopagebreak

\textbf{問}：$Q=\frac34$。走 4 步之後，變異數比普通醉漢的 4 大還是小？

\textbf{兩個狀態}：$L(N;x,y)$、$R(N;x,y)$＝走了 $N$ 步停在 $y$、最後一步向左／向右的機率。最後一步向左，上一步一定停在 $y+l$：
\begin{gather*}
L(N+1;x,y)=Q\,L(N;x,y+l)+(1-Q)\,R(N;x,y+l),\tag{5.11}\\
R(N+1;x,y)=(1-Q)\,L(N;x,y-l)+Q\,R(N;x,y-l).\tag{5.12}
\end{gather*}
\textbf{算到底}（$Q=\frac34$，起點 $L=R=\frac12$）：$N=1$ 是 $x=-1{:}\ (L,R)=(\frac12,0)$、$x=1{:}\ (0,\frac12)$。
$N=2$：$x=1$ 的 $R=\frac12$ 照同方向（$\times\frac34$）到 $x=2$ 得 $R=\frac38$，回頭（$\times\frac14$）到 $x=0$ 得 $L=\frac18$；對稱地 $x=0$ 也有 $R=\frac18$。
逐步做下去得表 \ref{five-tab-lr}（和暴力列舉 $2^N$ 條路線一致）。$N=2$ 的變異數是 $2\times\frac38\times4=3$。
\textbf{揭曉}：$N=1..4$ 的變異數是 $1,\ 3,\ \frac{11}2,\ \frac{33}4$，比普通醉漢的 $1,2,3,4$ 大；$Q=\frac14$（老是回頭）是 $1,\ 1,\ \frac32,\ \frac74$。慣性讓分布變寬。

\begin{table}[htbp]
\centering
\FiveTabLR\\[5pt]
{\small\FiveTabPVar}
\caption{有慣性的漫步。上：$Q=\frac34$ 的逐步表，每格是 $(L,R)$，兩數相加是停在 $x$ 的機率。
下左：小 $N$ 的變異數與持續長度 $n$（$Q=\frac12$ 是普通醉漢）；下右：大 $N$ 時精確值 $\to NQ/(1-Q)+$ 常數（5.2.2 節）}\label{five-tab-lr}\label{five-tab-pvar}
\end{table}

\textbf{轉移矩陣}：對 $y$ 取傅立葉，平移變成相位 $e^{\pm iql}$，$N$ 步就是 $\mathbf T^N$：
\begin{gather*}
\begin{pmatrix}L(N+1;q)\\R(N+1;q)\end{pmatrix}=\mathbf T\begin{pmatrix}L(N;q)\\R(N;q)\end{pmatrix},\qquad
\mathbf T=\begin{pmatrix}Qe^{-iql}&(1-Q)e^{-iql}\\(1-Q)e^{iql}&Qe^{iql}\end{pmatrix},\tag{5.15}\\
\lambda_{1,2}=Q\cos ql\pm\sqrt{(1-Q)^2-Q^2\sin^2ql}.\tag{5.17, 5.18}
\end{gather*}
（$\operatorname{tr}\mathbf T=2Q\cos ql$、$\det\mathbf T=2Q-1$。）$Q=\frac34$、$q=0.1$：$\lambda_1=0.9848$、$\lambda_2=0.5077$；$q=0.5$ 時 $Q|\sin q|>1-Q$，
兩個特徵值成共軛複數 $0.6582\pm0.2584i$。總數的傅立葉 $\sum_xP(N;x)e^{iqx}$（$q=0.3$）：$N=5$ 是 0.5699、$N=20$ 是 0.0388（三種算法一致）。

\textbf{生成函數}（(5.19)--(5.22)）：$\sum_Nz^N\mathbf T^N=(1-z\mathbf T)^{-1}$＝伴隨矩陣 $/D$，
\[D=1-2zQ\cos ql+z^2(2Q-1)=(1-z\lambda_1)(1-z\lambda_2),\qquad c_N\equiv[z^N]\frac1D=\frac{\lambda_1^{N+1}-\lambda_2^{N+1}}{\lambda_1-\lambda_2}.\tag{5.22, 5.25}\]
從 $L=R=\frac12$ 出發，總數的生成函數是 $\frac{1-z(2Q-1)\cos ql}{D}$，所以總數的傅立葉 $=c_N-(2Q-1)\cos ql\,c_{N-1}$。$N$ 大時由最大的 $\lambda_1$ 主導。

\subsubsection{持續長度與二階遞迴}
\leavevmode\bk{pp.~100--103}\ \anim{W503}\par\nopagebreak

\textbf{問}：$Q=\frac34$ 的醉漢走 1000 步，寬度相當於普通醉漢走幾步？

小 $q$ 展開最大的特徵值，$\lambda_1^N\approx e^{-NQq^2l^2/2(1-Q)}$，反傅立葉又是高斯：
\begin{gather*}
\lambda_1=1-\frac{Q}{1-Q}\frac{q^2l^2}{2}+O(q^4)\qquad(q=0.1{:}\ 0.985,\ \text{精確 }0.9848),\tag{5.23}\\
P(N;x)\approx\sqrt{\frac{1-Q}{2\pi NQl^2}}\exp\Big[-\frac{x^2(1-Q)}{2NQl^2}\Big].\tag{5.27}
\end{gather*}
變異數 $N\frac{Q}{1-Q}l^2$。精確值（$r=2Q-1$ 是相鄰兩步的相關，相隔 $k$ 步是 $r^k$）：
\[\langle x^2\rangle_N=l^2\Big[N\frac{1+r}{1-r}-\frac{2r(1-r^N)}{(1-r)^2}\Big]\ \xrightarrow{N\to\infty}\ N\frac{Q}{1-Q}l^2+\text{常數}.\]
表 \ref{five-tab-pvar}：精確值與 $NQ/(1-Q)$ 只差一個常數（$Q=\frac34$ 差 $-4$、$Q=\frac14$ 差 $+\frac49$）。$N=100$、$x=0$：精確 0.04587，$2\times$(5.27) 是 0.04607（奇偶性的因子 2）。

\textbf{持續長度}（(5.28)--(5.31)）：(5.27) 就是普通漫步的高斯，只是步長與步數換掉了：
\[l\to nl,\qquad N\to N/n,\qquad n=\frac{Q}{1-Q},\qquad l_p=nl\ \ \text{（持續長度）}.\tag{5.28--5.30}\]
$Q=\frac34,\ \frac12,\ \frac14$ 的 $n=3,\ 1,\ \frac13$。\textbf{揭曉}：$Q=\frac34$ 走 1000 步，變異數 $1000\times3=3000$（精確 2996），
相當於普通醉漢走 3000 步，也相當於步長 3 的醉漢走 333 步。

\textbf{兩個極端}：$Q=1$ 一直往同一邊走，變異數 $N^2$（$N=10$：100），是直線、永遠不是高斯（$n=\infty$）；
$Q=0$ 每步都回頭，在 0 與 $\pm1$ 之間來回：$N=10$ 的變異數 0、$N=11$ 的變異數 1。

\textbf{二階遞迴}（書 5.2.1 節）：從 (5.11)(5.12) 消去 $L$（(5.34)），只剩 $R$ 的遞迴（程式用逐步遞迴核對，殘差為 0）：
\[R(N+2;q)=2Q\cos ql\,R(N+1;q)+(1-2Q)\,R(N;q).\tag{5.35}\]
代入 $R=\lambda^N$ 得特徵方程 $\lambda^2-2Q\cos ql\,\lambda-(1-2Q)=0$（(5.37)），兩個根正是 (5.17)(5.18) 的 $\lambda_{1,2}$（根的乘積 $2Q-1=\det\mathbf T$）。
\begin{errbox}[p.~103，(5.37)]
書上印成 $\lambda^2-2Q\cos ql\,\lambda+(1-2Q)=0$。矩陣 (5.36) 的行列式是 $-(1-2Q)$，直接把 $R=\lambda^N$ 代入 (5.35) 也是 $-(1-2Q)$。
把 $\lambda_1=0.9848$（$q=0.1$、$Q=\frac34$）代入書上的式子，差 1.000（$=2|1-2Q|$）。
\end{errbox}

\subsubsection{電報方程}
\leavevmode\bk{pp.~103--109}\ \anim{W504}\par\nopagebreak

\textbf{問}：走路人以速率 $v=1$ 前進，平均每單位時間反轉 $\gamma=1$ 次。$t=3$ 時有沒有人在 $x=3$ 之外？一次都沒轉身的有多少？

\textbf{從格子到連續}（練習 5.3）：時間步 $\delta t$、步長 $l=v\delta t$、每步回頭的機率 $1-Q=\gamma\delta t$。$\delta t\to0$ 時 (5.11)(5.12) 變成
\[\frac{\partial R}{\partial t}+v\frac{\partial R}{\partial x}=\gamma(L-R),\qquad \frac{\partial L}{\partial t}-v\frac{\partial L}{\partial x}=\gamma(R-L):\tag{5.43, 5.44}\]
右行者隨波前進，並以速率 $\gamma$ 變成左行者，反之亦然。令 $\Sigma=R+L$（密度）、$\Delta=R-L$（流量 $\div v$），消去 $\Delta$（練習 5.4）：
\[\partial_t\Sigma+v\,\partial_x\Delta=0,\quad \partial_t\Delta+v\,\partial_x\Sigma=-2\gamma\Delta\quad\Longrightarrow\quad
\frac{\partial^2\Sigma}{\partial t^2}+2\gamma\frac{\partial\Sigma}{\partial t}=v^2\frac{\partial^2\Sigma}{\partial x^2}.\tag{5.45--5.48}\]
這就是\textbf{電報方程}：波動方程加上阻尼。格子核對（$t=1$、$x=0$，密度＝機率 $\div2l$）：$\delta t=0.01$ 得 0.3392、$\delta t=0.005$ 得 0.3380，
外插 $2\times0.3380-0.3392=0.3368$，正是下面公式的值。

\textbf{解}：起點 $\Sigma(x,0)=\delta(x)$、$\Delta=0$。取傅立葉與 Laplace（$t\to\lambda$）：
\[\big[\lambda^2+2\gamma\lambda+(qv)^2\big]\sigma(q,\lambda)=\lambda+2\gamma,\qquad
\sigma(q,t)=e^{-\gamma t}\Big[\cosh\omega t+\frac{\gamma}{\omega}\sinh\omega t\Big],\ \ \omega=\sqrt{\gamma^2-q^2v^2}.\tag{5.52}\]
反傅立葉（書上 (5.53)--(5.60) 用 $I_0$ 的積分表示），更正後是
\[\Sigma(x,t)=\frac{\gamma e^{-\gamma t}}{2v}\Big[I_0(\gamma s)+\frac{t}{s}I_1(\gamma s)\Big]\Theta(vt-|x|)+\frac{e^{-\gamma t}}{2}\big[\delta(x-vt)+\delta(x+vt)\big],\quad s=\sqrt{t^2-\frac{x^2}{v^2}}.\tag{5.61}\]
讀法：\textbf{兩個尖峰}是一次都沒轉身的人（各 $\frac12e^{-\gamma t}$），以速率 $v$ 往外跑；\textbf{連續部分}是轉過身的人，全在 $|x|<vt$ 之內——沒有人跑得比 $v$ 快（因果性）。
\begin{errbox}[p.~107，(5.57) 前、(5.60)(5.61)]
書上的前因子是 $\frac12\big(2\gamma+\frac{\partial}{\partial t}\big)$，少了 $1/v$（(5.58) 有）：應為 $\frac{1}{2v}\big(2\gamma+\frac{\partial}{\partial t}\big)\big[e^{-\gamma t}I_0(\gamma s)\big]$。
照書上算，$v=2$、$t=3$ 的總數是 1.9502，不是 1。(5.57) 前的「when $x<t$」應為 $x<vt$。
〔注〕$\partial_t$ 不作用在 $\Theta$（否則尖峰算兩次）；$\Theta(vt-x)$ 應為 $\Theta(vt-|x|)$；p.~106「Substituting into (5.50)」應為 (5.49)。
\end{errbox}

\textbf{算到底}（$v=\gamma=1$）：$t=1$、$x=0$ 時 $s=1$，$\Sigma=\frac{e^{-1}}2[I_0(1)+I_1(1)]=0.1839\times(1.2661+0.5652)=0.3368$。
前緣後方 $s\to0$：$I_0\to1$、$\frac tsI_1(\gamma s)\to\frac{\gamma t}2$，所以 $\Sigma(vt^-,t)=\frac{\gamma e^{-\gamma t}}{2v}(1+\frac{\gamma t}2)$，$t=1$ 是 0.2759。
守恆：連續部分的總量是 $1-e^{-\gamma t}$，加上兩個尖峰正好是 1（圖 \ref{five-fig-tele}、表 \ref{five-tab-tele}）。

\begin{figure}[htbp]
\centering
\FiveFigTele
\caption{書上 Fig.~5.3 的樣子（$v=\gamma=1$，只畫 $x\ge0$）：$t=1,3,5,8$ 的連續部分，箭頭是 $x=vt$ 的尖峰，標的是重量 $\frac12e^{-t}$。
灰虛線是 $t=8$ 的高斯（$D=\frac{v^2}{2\gamma}=\frac12$）。前緣以速率 $v$ 往外跑，尖峰越來越小，後面的人慢慢變成高斯}\label{five-fig-tele}
\end{figure}

\begin{table}[htbp]
\centering
\small
\FiveTabTele
\caption{電報方程的數字。$\Sigma(0,t)$ 從下面趨近高斯；變異數 $\to2Dt-\frac12$}\label{five-tab-tele}
\end{table}

\textbf{揭曉}：$t=3$ 時 $|x|>3$ 沒有人；沒轉過身的 $e^{-3}=0.0498$（在 $x=\pm3$ 的尖峰上；10 萬人的模擬 0.0488，$|x|<1.5$ 的比例 0.6060，公式 0.6051）。

\textbf{大 $t$}（練習 5.5--5.6）：$I_0(z)\approx I_1(z)\approx e^z/\sqrt{2\pi z}$、$s\approx t-\frac{x^2}{2v^2t}$，尖峰消失，剩下擴散係數 $D=\frac{v^2}{2\gamma}=0.5$ 的高斯。精確變異數
\[\langle x^2\rangle=\frac{v^2}{\gamma}\Big[t-\frac{1-e^{-2\gamma t}}{2\gamma}\Big]\ \xrightarrow{\ t\to\infty\ }\ 2Dt-\frac{v^2}{2\gamma^2}\]
（$t=1,3,5,8$：0.5677、2.5012、4.5000、7.5000）。持續長度 $l_p=v/\gamma=1$（平均跑多遠才轉身），$D=\frac12vl_p$：
和 5.2.2 節同一個道理（$Q=1-\gamma\delta t$、$l=v\delta t$ 代入 $N\frac{Q}{1-Q}l^2$ 正是 $\frac{v^2t}{\gamma}$）——慣性把步長放大成 $l_p$。

\subsubsection{二維、三維：跑一跑、轉個彎}
\leavevmode\bk{pp.~109--118}\ \anim{W505}\par\nopagebreak

\textbf{問}：大腸桿菌以速率 $v$ 游，平均每單位時間翻滾 $\gamma$ 次、換一個完全隨機的方向（$v=\gamma=1$）。$t=3$ 時有多少還在第一條直線上？前緣長什麼樣子？

\textbf{方程}：$f(\vec r,\hat n,t)$＝在 $\vec r$、朝 $\hat n$ 游的密度（二維 $\hat n=(\cos\theta,\sin\theta)$）：
\[\frac{\partial f}{\partial t}+v\cos\theta\frac{\partial f}{\partial x}+v\sin\theta\frac{\partial f}{\partial y}=-\Gamma f,\qquad \Gamma=\gamma-\gamma|0\rangle\langle0|.\tag{5.65, 5.69}\]
$|0\rangle$＝各方向均勻：每個方向以速率 $\gamma$ 流失、再平均分給所有方向（一維版就是 (5.43)(5.44)）。傅立葉＋Laplace 之後
$A=\lambda+iv\vec q\cdot\hat n+\gamma$ 對每個方向只是一個數，$\gamma|0\rangle\langle0|$ 是秩 1，反算子可以直接寫（練習 5.7）：
\[(A-\gamma|0\rangle\langle0|)^{-1}=A^{-1}+\frac{\gamma A^{-1}|0\rangle\langle0|A^{-1}}{1-\gamma\langle0|A^{-1}|0\rangle}.\tag{5.73--5.75}\]
（程式把角度離散成 64 格核對，誤差 $<10^{-12}$。）密度的轉換是 $\frac{g}{1-\gamma g}$，只需要角度平均 $g=\langle0|A^{-1}|0\rangle$
（$\lambda=q=1$ 時是 0.4472、0.4636；二維代入得 (5.78) 的 $\frac{1}{\sqrt{(\lambda+\gamma)^2+q^2v^2}-\gamma}$）：
\[g_{2D}=\frac{1}{\sqrt{(\lambda+\gamma)^2+q^2v^2}},\qquad g_{3D}=\frac{1}{vq}\arctan\frac{vq}{\lambda+\gamma}.\tag{5.77, 5.93}\]

\textbf{二維的答案}：反 Laplace 有極點與切割兩部分，書上 (5.84)、(5.89) 各算一部分（要相加）。極點給 $\frac{\gamma}{\pi v}e^{-\gamma t}\frac{\cosh(\gamma s/v)}{s}$、
切割給 $-\frac{\gamma}{2\pi v}e^{-\gamma t}\frac{e^{-\gamma s/v}}{s}$（$s=\sqrt{v^2t^2-r^2}$），$2\cosh-e^{-}=e^{+}$，合起來
\[\rho_{2D}(r,t)=\frac{e^{-\gamma t}}{2\pi r}\delta(r-vt)+\frac{\gamma}{2\pi v}\frac{e^{-\gamma t+\gamma s/v}}{s}\Theta(vt-r).\]
第一項是沒翻滾過的人（半徑 $vt$ 的圓環，總量 $e^{-\gamma t}$）；第二項在前緣像 $1/\sqrt{vt-r}$ 發散，但可積。
$t=3$ 的徑向密度 $P_2=2\pi r\rho$：$r=1,2,2.5,2.9$ 是 0.2978、0.4167、0.3941、0.4052；圓環 0.0498、其餘 0.9502。
\begin{errbox}[p.~112--113，(5.83)(5.84)(5.88)(5.89)]
(5.83) 後的對應應為 $b\to i\gamma/v$（書上 $i\gamma$），而且 $J_{-1/2}(z)=\sqrt{2/\pi z}\cos z$ 把 $\sqrt{\pi/2}$ 抵掉：
(5.84) 的 $2\gamma\sqrt{\pi/2}\,e^{-\gamma t}\cosh(\gamma\sqrt{v^2t^2-r^2})$ 應為 $2\gamma\,e^{-\gamma t}\cosh(\gamma\sqrt{v^2t^2-r^2}/v)$。
(5.88) 的 $\int_0^\infty qJ_0(qx)J_0(qy)\,dq$ 是 $\frac1x\delta(x-y)$，不是 $\frac{2\pi}{x}\delta(x-y)$，所以 (5.89) 兩項前面的 $2\pi$ 都不該有。
照書上的係數積分，$t=3$ 得 1.2658，不是 1。
\end{errbox}

\textbf{三維的答案}（(5.95)--(5.106)，只列結果）：極點在 $\lambda=qv\cot(qv/\gamma)-\gamma$（$0<q<\gamma\pi/2v$），切割兩邊的函數值差在 $\pm\frac\pi2$。
密度分成三部分：沒翻滾過的球殼 $e^{-\gamma t}$（(5.98)）；極點 (5.101) 加上切割 (5.103) 的第一項；以及 (5.103) 的最後一項——\textbf{只翻滾一次的人}：
\[\rho_1=\frac{\gamma e^{-\gamma t}}{4\pi v^2rt}\ln\frac{vt+r}{vt-r},\qquad \int\rho_1\,d^3r=\gamma te^{-\gamma t}\qquad\Big(\text{用 }\int_0^1u\ln\frac{1+u}{1-u}\,du=1\Big).\]
前緣的對數發散就是它造成的。$r>vt$ 時 (5.101) 與 (5.103) 第一項正好抵銷（$t=3$、$r=3.5$：0.0113 與 $-0.0113$）：因果性。
書上的歸一化是二維 $2\pi\rho$、三維 $2\pi^2\rho$（所以「$r^2\Sigma$ 積分得 $\pi/2$」；Fig.~5.6 的縱軸 $r^2\Sigma=\frac\pi2P_3$）。
\begin{errbox}[p.~115、p.~117]
p.~115（(5.96) 前）：切割右邊 $\lambda=-\gamma+\epsilon+ixvq$ 上的函數是 $\frac1{vq}\big(\frac\pi2-\frac i2\ln\frac{1+x}{1-x}\big)$（書上 $\frac1{vq}(\frac\pi2-i\ln((1+x)(1-x)))$），左邊是 $-\frac\pi2$。
(5.106) 的指數應為 $e^{+\gamma ut\cot u}$（書上 $e^{-\gamma ut\cot u}$）：只有 $+$ 才「正好抵銷 (5.101)」。
\end{errbox}

\begin{figure}[htbp]
\centering
\FiveFigRT
\caption{書上 Fig.~5.6 重畫（縱軸改成標準歸一化的 $P_3=4\pi r^2\rho$）：$t=1,3,5,8$，畫到 $r=0.99\,vt$（前緣對數發散）；箭頭是沒翻滾過的球殼，標的是重量 $e^{-t}$}\label{five-fig-rt}
\end{figure}

\begin{table}[htbp]
\centering
\small
\FiveTabRT\qquad
\FiveTabDims
\caption{左：三維的徑向密度 $P_3$（$\int P_3\,dr=1-e^{-t}$）與球殼的重量 $e^{-t}$。右：$t=3$ 時三個維度的比較
（一維看 $|x|$；模擬 10 萬個、種子 5；最後一列是連續部分在前緣 $r\to vt$ 的行為）}\label{five-tab-rt}
\end{table}

\textbf{揭曉}（$t=3$）：沒翻滾的 $e^{-3}=0.0498$（球殼），翻一次的 $3e^{-3}=0.1494$，其餘 0.8009。
維度越高前緣越尖（一維有限、二維 $1/\sqrt{\ }$、三維 $\ln$），但都可積：沒有人跑得比 $v$ 快，也一個都不少。

% =====================================================================
\subsection{連續時間漫步}
\leavevmode\bk{pp.~118--125}\quad 前面都是每單位時間走一步。這一節讓走路人每落腳一次，都要等一段隨機的時間才走下一步（像等公車）。

\subsubsection{等一下再走：等待時間}
\leavevmode\bk{pp.~118--122}\ \anim{W506}\par\nopagebreak

\textbf{問}：$t=3$ 之前走了幾步？「固定速率」和「越等越可能來」哪個走得多？

\textbf{等待}：$q(t)\,dt$＝已經等了 $t$、在接下來 $dt$ 內出發的機率。還在等的機率 $a(\tau)$ 滿足 $a'=-qa$：
\[a(\tau)=\exp\Big[-\int_0^\tau q(t)\,dt\Big],\qquad \int_0^\infty a(t)q(t)\,dt=1-a(\infty).\tag{5.108, 5.109}\]
$a(t)q(t)$ 是等待時間的密度。兩個例子（兩者的 $a(\infty)=0$：一定會走）：\textbf{固定速率} $q=p_0=1$：$a=e^{-t}$，$a(3)=e^{-3}=0.0498$；
\textbf{兩段等待}（像要轉兩班車，每段是速率 1 的指數等待）：等待密度 $aq=te^{-t}$、$a=(1+t)e^{-t}$、$q(t)=\frac{t}{1+t}$（越等越可能出發），$a(3)=4e^{-3}=0.1991$，平均等 2。

\textbf{走 $m$ 步＝卷積}：恰好走了一步、到 $\tau$ 還在等下一步的機率是 $\int_0^\tau a(t_1)q(t_1)a(\tau-t_1)\,dt_1$（(5.110)），$m$ 步是 $m$ 個 $aq$ 與一個 $a$ 的卷積（(5.111)）。
Laplace 轉換把卷積變乘法：
\begin{gather*}
A(\lambda)=\int_0^\infty aq\,e^{-\lambda t}dt,\qquad B(\lambda)=\int_0^\infty a\,e^{-\lambda t}dt,\qquad A=1-\lambda B,\tag{5.117--5.119}\\
\mathcal L[m\text{ 步}]=A^mB,\qquad \sum_mA^mB=\frac{B}{1-A}=\frac1\lambda.\tag{5.121, 5.122}
\end{gather*}
$A=1-\lambda B$ 來自分部積分（$aq=-a'$）；最後一式說機率加起來是 1（$\frac1\lambda$ 是常數 1 的 Laplace 轉換）。
$\lambda=0.7$：固定速率 $A=B=\frac1{1.7}=0.5882$；兩段等待 $A=\frac1{(1+\lambda)^2}=0.3460$、$B=\frac{2+\lambda}{(1+\lambda)^2}=0.9343$。

\textbf{算到底}：固定速率 $A^mB=\frac{p_0^m}{(p_0+\lambda)^{m+1}}$，反轉得 Poisson 分布 $P_m(t)=\frac{(p_0t)^m}{m!}e^{-p_0t}$；
兩段等待 $A^mB=\frac{1}{(1+\lambda)^{2m+1}}+\frac{1}{(1+\lambda)^{2m+2}}$，反轉得 $P_m(t)=e^{-t}\big[\frac{t^{2m}}{(2m)!}+\frac{t^{2m+1}}{(2m+1)!}\big]$
（走了 $m$ 步＝已經結束 $2m$ 或 $2m+1$ 段等待）。例：$P_1(3)=e^{-3}(\frac92+\frac{27}6)=9e^{-3}=0.4481$。
表 \ref{five-tab-ctrw}：平均步數 3 與 1.2506（$\frac t2-\frac14+\frac{e^{-2t}}4$）。\textbf{揭曉}：固定速率走得多——平均等 1 比平均等 2 短。

\subsubsection{等待＋步伐＝高斯；會消失的走路人}
\leavevmode\bk{pp.~122--125}\ \anim{W507}\par\nopagebreak

\textbf{問}：把等待和步伐一起算，走很久以後分布是什麼？時間變 4 倍，寬度變幾倍？

走了 $n$ 步的位移的傅立葉是 $\chi(\vec q)^n$（第 1 章），乘上「到 $t$ 恰好走了 $n$ 步」的 $A^nB$ 再加總：
\[\sum_n\chi(\vec q)^nA(\lambda)^nB(\lambda)=\frac{B(\lambda)}{1-\chi(\vec q)A(\lambda)}.\tag{5.124--5.128}\]
\textbf{練習 5.8}：固定速率時 $A=\frac{p_0}{p_0+\lambda}$、$B=\frac1{p_0+\lambda}$，(5.128) 變成 $\frac{1}{\lambda+p_0(1-\chi(q))}$；一維格子 $\chi=\cos q$，反轉得
$P(x,t)=e^{-p_0t}I_{|x|}(p_0t)$。$t=3$、$x=0,\pm1,\pm2$：0.2430、0.1968、0.1118。步數不固定，所以每一格都可能有人（沒有奇偶性）。

\textbf{高斯極限}：長距離、長時間（小 $q$、小 $\lambda$）：$\chi\approx1-Cq^2$、$A=1-\lambda B\approx1-\lambda B(0)$，
\begin{gather*}
\frac{B(0)}{1-(1-Cq^2)(1-\lambda B(0))}\approx\frac{1}{\lambda+(C/B(0))q^2}\ \Rightarrow\ e^{-Ctq^2/B(0)},\tag{5.129--5.131}\\
P(\vec r,t)\approx\Big(\frac{B(0)}{4\pi Ct}\Big)^{d/2}e^{-B(0)r^2/4Ct}.\tag{5.133 更正}
\end{gather*}
每個方向的變異數是 $2Ct/B(0)$。$B(0)=\int a\,dt$ 是平均等待時間：就是「每單位時間走 $1/B(0)$ 步」的普通漫步。
〔注〕(5.133) 本身是積分恆等式 $\int d^dq\,e^{i\vec q\cdot\vec r}e^{-Ctq^2/B(0)}=(\frac{\pi B(0)}{Ct})^{d/2}e^{-B(0)r^2/4Ct}$，機率還要乘 $(2\pi)^{-d}$；(5.131b) 的 $c$ 應為 $C$。
\begin{errbox}[p.~124，(5.133) 後]
書上說「width proportional to the time interval $t$」。應為\textbf{變異數} $\propto t$，寬度 $\propto\sqrt t$。
\end{errbox}
\textbf{算到底}（一維格子 $C=\frac12$，表 \ref{five-tab-ctrw}）：固定速率 $B(0)=1$，變異數 $t$：$t=10,40$ 是 10、40，寬度 3.162 $\to$ 6.325——
\textbf{揭曉：時間 $\times4$，寬度 $\times2$}。兩段等待 $B(0)=2$：精確變異數 4.75、19.75（$\approx t/2$；精確變異數就是平均步數）。
$P(0,10)$：固定速率精確 0.1278、高斯 0.1262；兩段等待 0.1805、0.1784。

\begin{table}[htbp]
\centering
\FiveTabCTRW
\caption{連續時間漫步。上：$t=3$ 走了 $m$ 步的機率；下：一維格子上的分布與高斯極限}\label{five-tab-ctrw}
\end{table}

\textbf{練習 5.9：會消失的走路人}：除了以速率 $p_0$ 走一步，還以速率 $p_1$ 消失。還在等的機率 $a=e^{-(p_0+p_1)t}$，但只有 $p_0$ 那部分是走一步：
$A=\frac{p_0}{p_0+p_1+\lambda}$、$B=\frac{1}{p_0+p_1+\lambda}$，所以 $\frac{B}{1-\chi A}=\frac{1}{\lambda+p_1+p_0(1-\chi)}$。$q=0$（$\chi=1$）就是總數 $\frac1{\lambda+p_1}$，即 $e^{-p_1t}$：衰減率是 $p_1$，和怎麼走無關。$p_0=1$、$p_1=0.1$、$t=10$：0.3679，半衰期 $\ln2/p_1=6.931$。〔注〕這時 $A=1-(\lambda+p_1)B\ne1-\lambda B$，(5.119) 不能照用。

% =====================================================================
\subsection{練習 5.1--5.9 解答}
{\small
\begin{enumerate}[label=\textbf{5.\arabic*},itemsep=0.05em]
\item 偏向漫步的擴散方程。\quad 答：多保留 $\frac12\partial_N^2P\approx\frac{v^2}2\partial_y^2P$，$D=\frac{l^2}2-\frac{v^2}2=2p(1-p)l^2$，只有 $p\to\frac12$ 時等於 $\frac{l^2}2$（5.1.1 節）。
\item 二階遞迴的解。\quad 答：$R=A_1\lambda_1^N+A_2\lambda_2^N$（$\lambda_{1,2}$＝更正後 (5.37) 的根），係數由 $R(0)=\frac12$、$R(1)=\frac12e^{iql}$ 定出；
加上 (5.34) 的 $L$，總數的傅立葉是 $c_N-(2Q-1)\cos ql\,c_{N-1}$（5.2.1 節）。
\item 由格子推 (5.43)(5.44)。\quad 答：$R(x,t+\delta t)=(1-\gamma\delta t)R(x-v\delta t,t)+\gamma\delta t\,L(x-v\delta t,t)$，展開到 $\delta t$ 一階：
$\delta t(\partial_tR+v\partial_xR)=\gamma\delta t(L-R)$，就是 (5.43)；$L$ 同理。
\item 電報方程。\quad 答：$\partial_t\Sigma=-v\partial_x\Delta$ 再對 $t$ 微分，代入 $\partial_t\Delta=-v\partial_x\Sigma-2\gamma\Delta$ 與 $v\partial_x\Delta=-\partial_t\Sigma$，
得 $\Sigma_{tt}+2\gamma\Sigma_t=v^2\Sigma_{xx}$；$\Delta$ 滿足同一個方程。
\item 大 $t$ 的高斯。\quad 答：$|x|\ll vt$ 時 $s\approx t-\frac{x^2}{2v^2t}$，$I_0(\gamma s)\approx I_1(\gamma s)\approx\frac{e^{\gamma s}}{\sqrt{2\pi\gamma t}}$，
$\Sigma\approx\frac{\gamma}{v}\frac{e^{-\gamma x^2/2v^2t}}{\sqrt{2\pi\gamma t}}=\frac{1}{\sqrt{4\pi Dt}}e^{-x^2/4Dt}$，$D=\frac{v^2}{2\gamma}$；尖峰 $e^{-\gamma t}$ 消失。$t=8$：0.1388 對 0.1410。
\item 變異數與持續長度。\quad 答：電報方程乘 $x^2$ 積分，$m(t)=\langle x^2\rangle$ 滿足 $\ddot m+2\gamma\dot m=2v^2$、$m(0)=\dot m(0)=0$，
所以 $m=\frac{v^2}{\gamma}\big[t-\frac{1-e^{-2\gamma t}}{2\gamma}\big]\approx\frac{v^2t}{\gamma}=2Dt$；持續長度 $l_p=v/\gamma$，$D=\frac12vl_p$。
\item 反算子 (5.75)。\quad 答：直接相乘（Sherman--Morrison），記 $g=\langle0|A^{-1}|0\rangle$：
$(A-\gamma|0\rangle\langle0|)A^{-1}=1-\gamma|0\rangle\langle0|A^{-1}$，而 $(A-\gamma|0\rangle\langle0|)\frac{\gamma A^{-1}|0\rangle\langle0|A^{-1}}{1-\gamma g}=\gamma|0\rangle\langle0|A^{-1}$，相加為 1。
\item 固定速率的連續時間漫步。\quad 答：(5.128) 得 $\frac1{\lambda+p_0(1-\chi(q))}$，反 Laplace 是 $e^{-p_0(1-\chi)t}$；一維格子 $\chi=\cos q$，
由 $e^{z\cos q}=\sum_nI_n(z)e^{inq}$ 得 $P(x,t)=e^{-p_0t}I_{|x|}(p_0t)$（5.3.2 節）。
\item 會消失的走路人。\quad 答：$\frac1{\lambda+p_1+p_0(1-\chi)}$，總數 $e^{-p_1t}$、衰減率 $p_1$（5.3.2 節；$t=10$：0.3679）。
\end{enumerate}}

% =====================================================================
\subsection{第 5 章的回顧與筆誤一覽}
\textbf{回顧}：改了規則，走路人最後還是回到高斯——只是平均、寬度與時間尺度換了樣子。
\begin{center}
\small\setlength{\tabcolsep}{4pt}
\begin{tabular}{@{}llll@{}}
\toprule
變化 & 結果 & 書上的式子 & 動畫\\
\midrule
偏向 & 移動的高斯，寬度 $\times\sqrt{4p(1-p)}$ & (5.1)、(5.7)、(5.10) & W501\\
慣性 & $\mathbf T$ 的特徵值；持續長度 $l_p=\frac{Q}{1-Q}l$ & (5.15)、(5.17)、(5.27)、(5.30) & W502--W503\\
電報 & 尖峰 $\frac12e^{-\gamma t}$＋$D=\frac{v^2}{2\gamma}$ 的高斯 & (5.48)、(5.61) & W504\\
二維、三維 & 前緣、因果性、守恆 & (5.69)、(5.75)、(5.101)、(5.103) & W505\\
等待 & $\frac{B}{1-\chi A}$，變異數 $2Ct/B(0)$ & (5.122)、(5.128)、(5.133) & W506--W507\\
\bottomrule
\end{tabular}
\end{center}
\textbf{筆誤一覽}（給 2004 年版）：$\Theta(vt-x)$、p.~106 的 (5.50)、(5.131b) 的 $c$、(5.133) 的 $(2\pi)^{-d}$、練習 5.9 的 $A\ne1-\lambda B$ 是〔注〕，
其餘在文中用紅框；每一條都用程式核對過。
\begin{center}
\footnotesize
\FiveTabErrata
\end{center}
